%=============================================================================!
% 10.F  SOLUTIONS TO THE EXERCISES                                            !
%=============================================================================!
\unnumbered{Chapter~\ref{ch:code}: Building a molecular dynamics code}
\label{sec:md-solutions}

\solutionto{ex:md-surface}The fraction is below 1\% when $(1 - 2/n)^3 >
0.99$, that is $1 - 2/n > 0.99^{1/3} = 0.99665549$, or $n > 2/0.00334451 =
597.995$. So $n = 598$, at which the fraction is 0.0099999, just below
1\%, and $N = 598^3 = \num{2.14e8}$ atoms. The estimate $6/n$ would give
$n = 600$.

\solutionto{ex:md-fractional}$\vb{r} = s_1\vb{a} + s_2\vb{b} + s_3\vb{c}$
reads $as_1 - \frac12as_2 = \frac12a$, $\frac{\sqrt3}2as_2 =
\frac{\sqrt3}2a$ and $cs_3 = \frac14c$. The second gives $s_2 = 1$, then
the first $s_1 = 1$, and the third $s_3 = \frac14$. Wrapping removes the
whole numbers, leaving $(0, 0, \frac14)$, the position $(0, 0, \frac14c)$:
the atom is a copy of the one at fractional coordinates $(0, 0, \frac14)$,
the first of graphite's four (\cref{sec:md-start}).

\solutionto{ex:md-graphite}$\vb{b}\times\vb{c} = (\frac{\sqrt3}2ac,
\frac12ac, 0)$, of length $ac$, and $\vb{a}\cdot(\vb{b}\times\vb{c}) =
\frac{\sqrt3}2a^2c$. Likewise $\vb{c}\times\vb{a} = (0, ac, 0)$, of length
$ac$, and $\vb{a}\times\vb{b} = (0, 0, \frac{\sqrt3}2a^2)$. The widths
$V/\lvert\vb{b}\times\vb{c}\rvert$, $V/\lvert\vb{c}\times\vb{a}\rvert$ and
$V/\lvert\vb{a}\times\vb{b}\rvert$ are $\frac{\sqrt3}2a$, $\frac{\sqrt3}2a$
and $c$: $V = \SI{35.29}{\angstrom\cubed}$, and the widths are 2.134,
2.134 and \SI{6.711}{\angstrom}.

\solutionto{ex:md-primitive}With the factor $\frac12a_0$ taken out of each
vector, $(1, 0, 1)\times(1, 1, 0) = (-1, 1, 1)$, of length $\sqrt3$, and
$(0, 1, 1)\cdot(-1, 1, 1) = 2$, so $V = 2(\frac12a_0)^3 = \frac14a_0^3$
and the width is $2(\frac12a_0)^3/(\sqrt3(\frac12a_0)^2) = a_0/\sqrt3$,
the same in all three directions by symmetry. For $a_0 =
\SI{5.267}{\angstrom}$: $V = \SI{36.53}{\angstrom\cubed}$, widths
\SI{3.041}{\angstrom}, and a cutoff of at most \SI{1.520}{\angstrom}.

\solutionto{ex:md-round}$\vb{s} = (0.9, -0.967, 0.025)$, which rounds to
$(1, -1, 0)$, so the copy is $\vb{d} - (2, -3, 0) = (-0.2, 0.1, 0.1)$\,\AA.
Each component can be changed only by a whole side, and $-0.2$, 0.1 and
0.1 are already within half a side of zero, so no copy is shorter.

\solutionto{ex:md-normal}With $u = (x - x_0)/(s\sqrt2)$, $\dd x =
s\sqrt2\,\dd u$ and the limits stay infinite, so the integral is
$s\sqrt2\int\mathrm{e}^{-u^2}\dd u = s\sqrt2\sqrt\pi = s\sqrt{2\pi}$.

\solutionto{ex:md-erf}$\mathrm{e}^{-u^2} = 1 - u^2 + \frac12u^4 - \dotsb$,
from the series of $\mathrm{e}^y$ (\cref{sec:ne-drag}) with $y = -u^2$. Integrating term by
term from 0 to $x$ gives $x - x^3/3 + x^5/10 - \dotsb$, so
$\operatorname{erf}(\alpha r)/r = (2/\sqrt\pi)(\alpha - \alpha^3r^2/3 +
\dotsb)$, which tends to $2\alpha/\sqrt\pi$ as $r \to 0$.

\solutionto{ex:md-orthogonal}The product is $\frac12[\cos(2\pi(m -
m')x/L) + \cos(2\pi(m + m')x/L)]$. For a whole number $k \neq 0$,
$\cos(2\pi kx/L)$ has the antiderivative $L\sin(2\pi kx/L)/(2\pi k)$,
which takes the same value at $x = 0$ and $x = L$, so its average over a
period is zero; for $k = 0$ the cosine is 1. If $m = m'$, then $m + m'
\neq 0$ and the average is $\frac12(1 + 0)$; if $m \neq m'$, neither $m -
m'$ nor $m + m'$ is zero and the average is 0.

\solutionto{ex:md-reciprocal}With $\vb{a} = (a, 0, 0)$, $\vb{b} =
(-\frac12a, \frac{\sqrt3}2a, 0)$ and $\vb{c} = (0, 0, c)$: $\vb{a}\cdot
\vb{a}^* = 2\pi$, $\vb{b}\cdot\vb{a}^* = \frac{2\pi}{a}(-\frac12a +
\frac{\sqrt3}2a\cdot\frac1{\sqrt3}) = 0$, $\vb{a}\cdot\vb{b}^* = 0$,
$\vb{b}\cdot\vb{b}^* = \frac{2\pi}{a}\cdot\frac{\sqrt3}2a\cdot
\frac2{\sqrt3} = 2\pi$, and $\vb{c}^*$ is at right angles to $\vb{a}$
and $\vb{b}$ with $\vb{c}\cdot\vb{c}^* = 2\pi$, while $\vb{a}^*$ and
$\vb{b}^*$ have no $z$ component. The lengths are $\lvert\vb{a}^*\rvert =
\lvert\vb{b}^*\rvert = \frac{2\pi}{a}\sqrt{1 + \frac13} = 4\pi/(\sqrt3a) =
\SI{2.9445}{\per\angstrom}$ and $\lvert\vb{c}^*\rvert = 2\pi/c =
\SI{0.9363}{\per\angstrom}$.

\solutionto{ex:md-forces}Only $\mathcal{C}$ and $\mathcal{S}$ depend on $\vb{r}_i$, through
the terms $q_i\cos(\vb{G}\cdot\vb{r}_i)$ and $q_i\sin(\vb{G}\cdot
\vb{r}_i)$, whose gradients (\cref{sec:en-gradient}) are $-q_i\sin(\vb{G}
\cdot\vb{r}_i)\,\vb{G}$ and $q_i\cos(\vb{G}\cdot\vb{r}_i)\,\vb{G}$. The
gradient of $\mathcal{C}^2 + \mathcal{S}^2$ is $2\mathcal{C}$ times the first plus
$2\mathcal{S}$ times the
second, $2q_i\vb{G}\,[\mathcal{S}(\vb{G})\cos(\vb{G}\cdot\vb{r}_i) -
\mathcal{C}(\vb{G})\sin(\vb{G}\cdot\vb{r}_i)]$, and $\vb{F}_i$ is minus
$2\pi k_{\mathrm e}/V$ times the sum of $\mathrm{e}^{-G^2/(4\alpha^2)}/G^2$
times that; the minus sign reverses the bracket and the 2 makes $4\pi$,
which is the formula. For
the real-space term, $\frac{\dd}{\dd r}\operatorname{erfc}(\alpha r) =
-(2/\sqrt\pi)\,\alpha\,\mathrm{e}^{-\alpha^2r^2}$ by the fundamental theorem
of calculus and the chain rule, and the product rule gives the slope of
$\operatorname{erfc}(\alpha r)\cdot r^{-1}$ as $-\operatorname{erfc}(\alpha
r)/r^2 - (2\alpha/\sqrt\pi)\,\mathrm{e}^{-\alpha^2r^2}/r$.

\solutionto{ex:md-rocksalt}The cell is $2\vb{I}$, so the reciprocal
vectors are $\pi$ along the axes and $\vb{G} = \pi(m_1, m_2, m_3)$. For a cation at
$\vb{p}$, $\vb{G}\cdot\vb{p} = \pi\,\vb{m}\cdot\vb{p}$ is a whole multiple
of $\pi$, and for its anion $\pi m_1$ more, so every sine is zero and
$\mathcal{S} = 0$. By the addition formula, $\cos(x + \pi m_1) = (-1)^{m_1}\cos x$,
so $\mathcal{C} = (1 - (-1)^{m_1})\sum_{\vb{p}}\cos(\pi\,\vb{m}\cdot\vb{p})$, and,
since $\cos k\pi = (-1)^k$ for a whole number $k$, the sum over the four
cations is $1 + (-1)^{m_2 + m_3} + (-1)^{m_1 + m_3} +
(-1)^{m_1 + m_2}$. If the $m_k$ are all even or all odd, every exponent is
even and the sum is 4; otherwise one of them differs in parity from the
other two, the two sums that contain it are odd and the third even, and
the sum is $1 + 1 - 1 - 1 = 0$. The factor $1 - (-1)^{m_1}$ is 2 for odd
$m_1$ and 0 for even, so $\mathcal{C} = 8$ when all are odd and 0 otherwise.

\solutionto{ex:md-cost}The copies within $r_{\mathrm c}$ number about
$\frac43\pi r_{\mathrm c}^3\rho \propto \rho/\alpha^3$; the waves within
$G_{\mathrm c}$ number about $\frac43\pi G_{\mathrm c}^3$ divided by the
volume per wave, $(2\pi)^3/V$, the cube of the spacing $2\pi/L$ for a cube
of side $L$, so $\propto V\alpha^3$. The cost is $T = C_1N\rho/\alpha^3 +
C_2NV\alpha^3$ for constants $C_1$ and $C_2$, and $\dd T/\dd\alpha =
-3C_1N\rho/\alpha^4 + 3C_2NV\alpha^2 = 0$ gives $\alpha^6 =
C_1\rho/(C_2V) = (C_1/C_2)\,N/V^2$. There the two terms are equal, and $T
= 2C_2NV\alpha^3 \propto NV(N/V^2)^{1/2} = N^{3/2}$.

\solutionto{ex:md-parameters}$\sqrt{-\ln10^{-6}} = 3.717$, so $\alpha =
3.717/\SI{15}{\angstrom} = \SI{0.2478}{\per\angstrom}$ and $G_{\mathrm c}
= 2\times0.2478\times3.717 = \SI{1.842}{\per\angstrom}$. The waves form a
cubic grid of spacing $2\pi/30 = \SI{0.2094}{\per\angstrom}$, one per
$(2\pi/30)^3$ of volume, so the sphere of radius $G_{\mathrm c}$ holds
about $\frac43\pi(1.842)^3/(0.2094)^3 = 2850$; counting the grid points
gives 2896.

\solutionto{ex:md-water}$\frac43\pi(10)^3\times0.0334 = 139.9$, about
140 molecules.

\solutionto{ex:md-cells}$\lfloor40/9.5\rfloor = 4$, so $4^3 = 64$ small
cells; $\lfloor30/9.5\rfloor = 3$, 27 small cells; and $\lfloor25/9.5
\rfloor = 2$, fewer than three, so all pairs are tested.

\solutionto{ex:md-rebuild}$\sqrt{2K/m} = \sqrt{2\times0.0104/(39.948\times
103.64)} = \SI{0.00224}{\angstrom\per\femto\second}$, with the factor
\texttt{MV2\_TO\_EV} of \cref{eq:en-mv2}, or \SI{0.0224}{\angstrom}
a step, so \SI{0.5}{\angstrom} takes 22 steps. The run rebuilds every
10.7 steps because it rebuilds as soon as any one of 500 atoms has moved
half the skin, and some atoms move faster than the root-mean-square speed.

\solutionto{ex:md-drift}$\sum_im_i(\vb{v}_i - \vb{P}/M) = \vb{P} -
M\vb{P}/M = \vb{0}$, and $\sum_im_i(c\vb{v}_i) = c\sum_im_i\vb{v}_i =
\vb{0}$.

\solutionto{ex:md-half}Started at the minimum, $x = A\sin\omega t$ and
$K = \frac12m\omega^2A^2\cos^2\omega t$, with $K_0 = \frac12m\omega^2A^2$.
The average of $\cos^2$ over a period is $\frac12$ (\cref{sec:os-energy}),
so $\overline{K} = \frac12K_0$.

\solutionto{ex:md-femtosecond}The run in ASE's units lasts 1000 of its
femtoseconds, $1000\times\num{4.6e-9} = \SI{4.6e-6}{\femto\second}$
longer than \texttt{mdlab}'s, and starts with velocities smaller by the
same fraction. For an atom in free flight the two cancel; a vibrating
atom keeps its frequency whatever its speed, so the extra time carries it
on by about its speed times \SI{4.6e-6}{\femto\second}, $0.0031\times
\num{4.6e-6} = \num{1.4e-8}$\,\AA: the size of the difference found.
